\newcommand{\TabOptimisationParameters}{
\begin{table}[tb]
\begin{tabular}{ll}
	\hline
	Region for optimisation & Approx.\ No.\  of parameters \\
	\hline
	Production target dimensions and location & $3+3$ \\
	Torus1 dipole field strength & 1 \\
	Torus2 dipole field strength & 1 \\
	Muon beam collimator shapes, position, and material & $3+1+1$ \\
	Stopping target shape and location & $4+3$ \\
	Beam blocker position, form, and material& $3+3+1$ \\
	Electron spectrometer dipole field strength& 1 \\
	DIO blockers in the spectrometer & $4$ \\
	\hline
	Approx. total number of parameters& 32 \\
	\hline
\end{tabular}
	\caption{\tablabel{optimisation:possible-parameters}%
	The list of parameters that can be optimised and an estimate for the number of parameters that this represents.
	In the case of the target, beam blocker and collimator shapes the number of parameters is only approximate; crudely speaking there is at least a width, length and height but in principle one could have a very irregular shape that cannot be parametrised by only three numbers, for example shapes that change as a function of distance along the beamline.
}
\end{table}
}

\newcommand{\FigOptimProdTgtLength}{
\begin{figure}[t]
\centering
\subfloat[][\figlabel{optimisation:ProdTgtSec:Length:Momentum:Muons}Muons]{\includegraphics[width=0.45\textwidth,trim=0 0 0 1.5cm,clip]{figs/optimisation/ProdTgtGeom/Length_mu-minus_momentum}}
\subfloat[][\figlabel{optimisation:ProdTgtSec:Length:Momentum:Pions}Pions]{\includegraphics[width=0.45\textwidth,trim=0 0 0 1.5cm,clip]{figs/optimisation/ProdTgtGeom/Length_pi-minus_momentum}}
\caption{
Change to momentum distributions at the entrance to the first 90 degrees of the bent muon beam solenoid for different target lengths.
Target length is given as half-length which is the Geant4 convention.  
}
\label{optimisation:ProdTgtSec:Length:Momentum}
\end{figure}

\begin{figure}[t]
\centering
\subfloat[][\figlabel{optimisation:ProdTgtSec:Length:Integral:Muons}Muons]{\includegraphics[width=0.45\textwidth,trim=0 0 0 1.5cm,clip]{figs/optimisation/ProdTgtGeom/Length_mu-minus_integral_toZero}}
\subfloat[][\figlabel{optimisation:ProdTgtSec:Length:Integral:Pions}Pions]{\includegraphics[width=0.45\textwidth,trim=0 0 0 1.5cm,clip]{figs/optimisation/ProdTgtGeom/Length_pi-minus_integral_toZero}}
\caption{\figlabel{optimisation:ProdTgtSec:Length:Integral}
Integrated muon and pion yields up to a certain momentum at the entrance to the first 90 degrees of the bent muon beam solenoid as a function of target length.
}
\end{figure}
\begin{figure}[t]
\centering
\subfloat[][\figlabel{optimisation:ProdTgtSec:Length:IntegralRatio:Muons}Muons]{\includegraphics[width=0.45\textwidth,trim=0 0 0 1.5cm,clip]{figs/optimisation/ProdTgtGeom/Length_mu-minus_integral_ratios}}
\subfloat[][\figlabel{optimisation:ProdTgtSec:Length:IntegralRatio:Pions}Pions]{\includegraphics[width=0.45\textwidth,trim=0 0 0 1.5cm,clip]{figs/optimisation/ProdTgtGeom/Length_pi-minus_integral_ratios}}
\caption{\figlabel{optimisation:ProdTgtSec:Length:IntegralRatio}
Change in the momentum distribution of muons and pions at the entrance to the first 90 degrees of the bent muon beam solenoid as a function of target length.
}
\end{figure}
}

\newcommand{\FigOptimProdTgtRad}{
\begin{figure}[t]
\centering
\subfloat[][\figlabel{optimisation:ProdTgtSec:Radius:Momentum:Muons}Muons]{\includegraphics[width=0.45\textwidth,trim=0 0 0 1.5cm,clip]{figs/optimisation/ProdTgtGeom/Radius_mu-minus_momentum}}
\subfloat[][\figlabel{optimisation:ProdTgtSec:Radius:Momentum:Pions}Pions]{\includegraphics[width=0.45\textwidth,trim=0 0 0 1.5cm,clip]{figs/optimisation/ProdTgtGeom/Radius_pi-minus_momentum}}
\caption{
Change to momentum distributions at the entrance to the first 90 degrees of the bent muon beam solenoid for different target radii.
}
\label{optimisation:ProdTgtSec:Radius:Momentum}
\end{figure}

\begin{figure}[t]
\centering
\subfloat[][\figlabel{optimisation:ProdTgtSec:Radius:Integral:Muons}Muons]{\includegraphics[width=0.45\textwidth,trim=0 0 0 1.5cm,clip]{figs/optimisation/ProdTgtGeom/Radius_mu-minus_integral_toZero}}
\subfloat[][\figlabel{optimisation:ProdTgtSec:Radius:Integral:Pions}Pions]{\includegraphics[width=0.45\textwidth,trim=0 0 0 1.5cm,clip]{figs/optimisation/ProdTgtGeom/Radius_pi-minus_integral_toZero}}
\caption{\figlabel{optimisation:ProdTgtSec:Radius:Integral}
Integrated muon and pion yields up to a certain momentum at the entrance to the first 90 degrees of the bent muon beam solenoid as a function of target radius.
}
\end{figure}
\begin{figure}[t]
\centering
\subfloat[][\figlabel{optimisation:ProdTgtSec:Radius:IntegralRatio:Muons}Muons]{\includegraphics[width=0.45\textwidth,trim=0 0 0 1.5cm,clip]{figs/optimisation/ProdTgtGeom/Radius_mu-minus_integral_ratios}}
\subfloat[][\figlabel{optimisation:ProdTgtSec:Radius:IntegralRatio:Pions}Pions]{\includegraphics[width=0.45\textwidth,trim=0 0 0 1.5cm,clip]{figs/optimisation/ProdTgtGeom/Radius_pi-minus_integral_ratios}}
\caption{\figlabel{optimisation:ProdTgtSec:Radius:IntegralRatio}
Change in the momentum distribution of muons and pions at the entrance to the first 90 degrees of the bent muon beam solenoid as a function of target radius.
}
\end{figure}
}

\newcommand{\FigOptimMuBeamDipoleMuStops}{
\begin{figure}[b]
\centering
%	\fbox{
\includegraphics[width=0.85\textwidth,trim=0 0.5cm 0 1.0cm,clip]{figs/optimisation/MuonBeamDipoles/Tidied_stopped_muons.pdf}
%}
\caption{\figlabel{optim:muBeamDipole:stoppedMu}
	Muon stopping rate as a function of the two dipole field strengths (given relative to the \phaseI design specification).
	A clear anti-correlation is visible which is discussed in the text.
}
\end{figure}
}

\newcommand{\FigOptimMuBeamDipolePiStops}{
\begin{figure}[t]
\centering
%	\fbox{
\includegraphics[width=0.85\textwidth,trim=0 0.5cm 0 1.0cm,clip]{figs/optimisation/MuonBeamDipoles/Tidied_stopped_pions.pdf}
%}
\caption{\figlabel{optim:muBeamDipole:stoppedPi}
	Pion stopping rate as a function of the two dipole field strengths (given relative to the \phaseI design specification).
At the level of statistics used to generate each point, no clear trend is obvious.
Empty squares are those where no pions stopped in the run.
}
\end{figure}
}

\newcommand{\FigOptimMuBeamDipoleMuDispersion}{
\begin{figure}[b]
\centering 
\subfloat[][\figlabel{optim:MuBeamDipole:MuDispersion:Entry}Torus1 Entry]{\includegraphics[width=0.45\textwidth,trim=0 0.9cm 0 1.9cm,clip]{figs/optimisation/MuonBeamCollimators/Tidied_Muon_dispersion_at_entrance.pdf}}
\subfloat[][\figlabel{optim:MuBeamDipole:MuDispersion:Exit}Torus2 Exit]  {\includegraphics[width=0.45\textwidth,trim=0 0.9cm 0 1.9cm,clip]{figs/optimisation/MuonBeamCollimators/Tidied_Muon_dispersion_at_exit.pdf}}
\caption{\figlabel{optim:MuBeamDipole:MuDispersion}
Dispersive effect of the 180\degree bent transport solenoid and dipole field on muons.
No collimating material is yet included, so the high-energy muons being removed is due purely to the beam-pipe itself.
}
\end{figure}
}

\newcommand{\FigOptimMuBeamCollimMuonPaths}{
\begin{figure}[p]
\centering 
\subfloat[][\figlabel{optim:MuBeamCollim:Beamline:All}All Muons]                  {\includegraphics[width=1\textwidth,trim=6cm 0.5cm 2.5cm 1.9cm,clip]{figs/optimisation/MuonBeamCollimators/Tidied_All_muons_wGeom}}\\
\subfloat[][\figlabel{optim:MuBeamCollim:Beamline:Stopped}Stopped Muons]          {\includegraphics[width=1\textwidth,trim=6cm 0.5cm 2.5cm 1.9cm,clip]{figs/optimisation/MuonBeamCollimators/Tidied_Stopped_muons_wGeom}}\\
\subfloat[][\figlabel{optim:MuBeamCollim:Beamline:HighP}Muons with $p>70$ MeV/c around the stopping target]          {\includegraphics[width=1\textwidth,trim=6cm 0.5cm 2.5cm 1.9cm,clip]{figs/optimisation/MuonBeamCollimators/Tidied_HighP_muons_wGeom}}\\
\subfloat[][\figlabel{optim:MuBeamCollim:Beamline:Diff}Where to Place Collimators]{\includegraphics[width=1\textwidth,trim=5cm 0.5cm 1.5cm 1.3cm,clip]{figs/optimisation/MuonBeamCollimators/Tidied_Where_to_collimate_wGeom}}
\caption{\figlabel{optim:MuBeamCollim:Beamline}
The heights of muons as they pass along the beamline.  
	\protect\subref{fig:optim:MuBeamCollim:Beamline:All} The path of all muons.
	\protect\subref{fig:optim:MuBeamCollim:Beamline:Stopped}: The paths of muons that stop in the target.
	\protect\subref{fig:optim:MuBeamCollim:Beamline:HighP}: The heights of muons with momentum greater than 70 MeV/c when they enter the region around the stopping target.  These could potentially decay in flight to give electrons with 100 MeV/c or greater.
	\protect\subref{fig:optim:MuBeamCollim:Beamline:Diff}: The difference between plot \protect\subref{fig:optim:MuBeamCollim:Beamline:Stopped} and plot \protect\subref{fig:optim:MuBeamCollim:Beamline:HighP}.  Regions in dark blue would give the greatest impact in removing high momentum muons whilst leave the stopping muons untouched.
	These plots should be compared to those of \fig{optim:MuBeamCollim:BeamWColl} once collimators have been introduced.
}
\end{figure}
}

\newcommand{\FigOptimMuBeamCollimTransverseSep}{
\begin{figure}[t]
\centering 
\subfloat[][\figlabel{optim:MuBeamCollim:TransverseSep:start}At 0\degree]  {\includegraphics[height=0.35\textheight,trim=0.0cm 0.8cm 1.3cm 1.9cm,clip]{figs/optimisation/MuonBeamCollimators/MuonTransversePos_Torus1Tor1FirstPoint}}
\subfloat[][\figlabel{optim:MuBeamCollim:TransverseSep:start}At 90\degree] {\includegraphics[height=0.35\textheight,trim=1.7cm 0.8cm 1.3cm 1.9cm,clip]{figs/optimisation/MuonBeamCollimators/MuonTransversePos_Torus2Monitor_0}}
\subfloat[][\figlabel{optim:MuBeamCollim:TransverseSep:start}At 180\degree]{\includegraphics[height=0.35\textheight,trim=1.7cm 0.8cm 1.3cm 1.9cm,clip]{figs/optimisation/MuonBeamCollimators/MuonTransversePos_StopTgtSecMonitor_0}}
\caption{\figlabel{optim:MuBeamCollim:TransverseSep}
The separation between stopping and dangerous muons.
The separation is largest at the exit (180\degree), reasonable at the entrance (0\degree), and smallest around the mid-point (90\degree).
}
\end{figure}
}

\newcommand{\FigOptimMuBeamCollimMuonPathsWColl}{
\begin{figure}[p!]
\centering 
\subfloat[][\figlabel{optim:MuBeamCollim:BeamWColl:All}All Muons]                  {\includegraphics[width=1\textwidth,trim=6cm 0.5cm 2.5cm 1.9cm,clip]{figs/optimisation/MuonBeamCollimators/Tidied_WColl_AllMuons_WGeom}}\\
\subfloat[][\figlabel{optim:MuBeamCollim:BeamWColl:Stopped}Stopped Muons]          {\includegraphics[width=1\textwidth,trim=6cm 0.5cm 2.5cm 1.9cm,clip]{figs/optimisation/MuonBeamCollimators/Tidied_WColl_StoppedMuons_WGeom}}\\
\subfloat[][\figlabel{optim:MuBeamCollim:BeamWColl:HighP}Muons with $p>70$ MeV/c around the stopping target]          {\includegraphics[width=1\textwidth,trim=6cm 0.5cm 2.5cm 1.9cm,clip]{figs/optimisation/MuonBeamCollimators/Tidied_WColl_HighPMuons_WGeom}}\\
\caption{\figlabel{optim:MuBeamCollim:BeamWColl}
The heights of muons as they pass along the beamline.  
	\protect\subref{fig:optim:MuBeamCollim:Beamline:All} The path of all muons.
	\protect\subref{fig:optim:MuBeamCollim:Beamline:Stopped}: The paths of muons that stop in the target.
	\protect\subref{fig:optim:MuBeamCollim:Beamline:HighP}: The heights of muons with momentum greater than 70 MeV/c when they enter the region around the stopping target.  These could potentially decay in flight to give electrons with 100 MeV/c or greater.
	These plots should be compared to those of \fig{optim:MuBeamCollim:Beamline} before collimators were introduced, where it is clear how well the dangerous muons are being suppressed.
}
\end{figure}
}

\newcommand{\FigOptimMuBeamCollimTorusOne}{
\begin{figure}[t]
\centering 
\subfloat[][\figlabel{optim:MuBeamCollim:Torus1:perPOT}Particles Surviving per POT]{\includegraphics[width=0.5\textwidth,trim=0.8cm 0.8cm 0.6cm 1.9cm,clip]{figs/optimisation/MuonBeamCollimators/Survived_Coll1_unNormalised-log.pdf}}
\subfloat[][\figlabel{optim:MuBeamCollim:Torus1:fraction}Fraction Surviving Collimator]{\includegraphics[width=0.5\textwidth,trim=0.8cm 0.8cm 0.6cm 1.9cm,clip]{figs/optimisation/MuonBeamCollimators/Survived_Coll1_Normalised-lin.pdf}}
\caption{\figlabel{optim:MuBeamCollim:Torus1}
The effect of changing the height of the collimator in Torus1 on the particle distributions.
}
\end{figure}
}

\newcommand{\FigOptimMuBeamCollimTorusTwoPerPOT}{
\begin{figure}[t]
\centering 
\includegraphics[width=0.95\textwidth,trim=0.3cm 0.3cm 0.8cm 0.1cm,clip]{figs/optimisation/MuonBeamCollimators/Survived_Coll2_unNormalised-lin.pdf}
\caption{\figlabel{optim:MuBeamCollim:Torus2:perPOT}
The number of particles reaching the end of the Torus2 solenoid per POT for different heights of both collimators in Torus1 and Torus2.
}
\end{figure}
}


\newcommand{\FigOptimMuBeamCollimTorusTwoFraction}{
\begin{figure}[t]
\centering 
\includegraphics[width=0.8\textwidth,trim=0.3cm 0.3cm 0.8cm 0.1cm,clip]{figs/optimisation/MuonBeamCollimators/Survived_Coll2_unNormalised-lin.pdf}
\caption{\figlabel{optim:MuBeamCollim:Torus2:fraction}
	The number of particles reaching the end of the Torus2 solenoid relative to the number that enter the Torus1 solenoid (\ie the survival probability) for different heights of both collimators in Torus1 and Torus2.
}
\end{figure}
}

\newcommand{\FigOptimMuBeamCollimTorusTwoContours}{
\begin{figure}[tb]
\centering 
\includegraphics[width=0.6\textwidth,trim=0.3cm 0.3cm 0.8cm 0.1cm,clip]{figs/optimisation/MuonBeamCollimators/Survived_Coll2_StoppedVsHighP-Muons.pdf}
\caption{\figlabel{optim:MuBeamCollim:Torus2:contours}
2.5\% contours showing how changing the collimator height affects the stopped and high-momentum muon rates.
Each contour represents a reduction of 2.5 percentage points to the particle flux, with 100\% acceptance sitting in the bottom right.
Blue lines represent the flux of stopping muons, whilst red lines represent the flux of high-momentum muons.
For example, for collimator heights within the first blue contour towards the bottom-right corner, less than 2.5\% of stopped muons are lost.
}
\end{figure}
}

\newcommand{\FigOptimESTDipoleBeamHeightTwoD}{
\begin{figure}[tp]
\centering 
	\subfloat[][\figlabel{optim:ESTDipole:Beam:0.0}No Dipole]  {\includegraphics[width=1\textwidth,trim=4cm 0.5cm 9.5cm 0.7cm,clip]{figs/optimisation/EST_dipole/Tidied_signal_height-dipole_00}}\\
\subfloat[][\figlabel{optim:ESTDipole:Beam:0.1}0.1~T Dipole]{\includegraphics[width=1\textwidth,trim=4cm 0.5cm 9.5cm 0.7cm,clip]{figs/optimisation/EST_dipole/Tidied_signal_height-dipole_10}}\\
\subfloat[][\figlabel{optim:ESTDipole:Beam:0.2}0.2~T Dipole]{\includegraphics[width=1\textwidth,trim=4cm 0.5cm 9.5cm 0.7cm,clip]{figs/optimisation/EST_dipole/Tidied_signal_height-dipole_20}}\\
\caption{\figlabel{optim:ESTDipole:Beam}
The heights of signal electrons for different dipole field values.
}
\end{figure}
}

\newcommand{\FigOptimESTDipoleBeamHeightMean}{
\begin{figure}[tb]
\centering 
%	\fbox{
\includegraphics[width=0.8\textwidth,trim=1.0cm 1.3cm 3.8cm 0.9cm,clip]{figs/optimisation/EST_dipole/Tidied_NoShift-Height}
%}
\caption{\figlabel{optim:ESTDipole:MeanHeight}
Mean height of signal electrons for different values of the dipole field strength.
}
\end{figure}
}

\newcommand{\FigOptimESTDipoleBeamFluxMean}{
\begin{figure}[tb]
\centering 
\includegraphics[width=0.8\textwidth,trim=1.0cm 0.3cm 3.8cm 1.8cm,clip]{figs/optimisation/EST_dipole/Tidied_NoShift-Flux}
\caption{\figlabel{optim:ESTDipole:MeanFlux}
Survival probability for signal electrons as a function of the distance along the beamline for different values of the electron spectrometer's dipole field strengths.
}
\end{figure}
}

\newcommand{\FigOptimESTDipoleAcceptanceVsDipole}{
\begin{figure}[tb]
\centering 
\includegraphics[width=0.6\textwidth,trim=0.3cm 0.3cm 1.9cm 1.1cm,clip]{figs/optimisation/EST_dipole/Tidied_acceptance}
\caption{\figlabel{optim:ESTDipole:acceptance}
Geometric acceptance into the StrECAL detector as a function of the dipole field strength over the electron spectrometer.
}
\end{figure}
}

\newcommand{\FigOptimStopTgtPosMuStops}{
\begin{figure}[bt]
\centering 
%	\fbox{
\includegraphics[width=0.7\textwidth,trim=2.0cm 0.05cm 0.9cm 0.3cm,clip]{figs/optimisation/StopTgtPosition/Tidied_MuonStoppingRate.pdf}
%}
\caption{\figlabel{optim:StopTgtPos:MuStops}
Muon stopping rate per \ac{POT} for different target positions.
The linear behaviour arises from the reduced field strength and fixed target radius such that fewer muons impact the target as it is moved downstream.
}
\end{figure}
}

\newcommand{\FigOptimStopTgtPosSensitivitySpect}{
\begin{figure}[tb]
\centering 
%	\fbox{
\includegraphics[width=0.9\textwidth,,trim=1.3cm 0.2cm 0.7cm 0.6cm,clip]{figs/optimisation/StopTgtPosition/Tidied_-AcceptedMomentum.pdf}
%}
\caption{\figlabel{optim:StopTgtPos:AcceptedMomSpect}
The momentum dependence of the electron acceptance into the detector for different target positions.
The spectrum for each target position is normalised to the muon stopping rate for that position, such that each 
The linear behaviour arises from the reduced field strength and fixed target radius such that fewer muons impact the target as it is moved downstream.
}
\end{figure}
}

\newcommand{\FigOptimStopTgtPosSensitivityNoBeamBlock}{
\begin{figure}[tb]
\centering 
%	\fbox{
\includegraphics[width=0.9\textwidth,trim=1.3cm 0.2cm 0.7cm 0.6cm,clip]{figs/optimisation/StopTgtPosition/Tidied_NoBeamBlocker-AcceptedMomentum.pdf}
%}
\caption{\figlabel{optim:StopTgtPos:AcceptedMomSpectNoBeamBlock}
The momentum dependence of the electron acceptance into the detector for different target positions when the beam blocker is removed.
}
\end{figure}
}

\newcommand{\FigOptimStopTgtPosSensitivityIntegral}{
\begin{figure}[tb]
\centering 
%	\fbox{
\subfloat[][\figlabel{optim:StopTgtPos:AcceptIntegral:Adjacent}Sensitivity]{%
\includegraphics[width=0.49\textwidth,trim=0.5cm 0.5cm 0.3cm 1.9cm,clip]{figs/optimisation/StopTgtPosition/Tidied_-AcceptedMomentum-Integrated_adjacent.pdf}}
\subfloat[][\figlabel{optim:StopTgtPos:AcceptIntegral:Ratio}Change in Shape]{%
\includegraphics[width=0.49\textwidth,trim=0.5cm 0.5cm 0.3cm 1.9cm,clip]{figs/optimisation/StopTgtPosition/Tidied_-AcceptedMomentum-Integrated_ratios.pdf}}
%}
\caption{\figlabel{optim:StopTgtPos:AcceptIntegral}
\protect\subref{fig:optim:StopTgtPos:AcceptIntegral:Adjacent} 
The variation in sensitivity (acceptance $\times$ stopping rate) to electrons with different momenta as a function of the target position with respect to the nominal location. 
The darkest red line towards the top of the plot represents the sensitivity to signal, and it is that line that should therefore be maximised.
\protect\subref{fig:optim:StopTgtPos:AcceptIntegral:Ratio} 
The change in the shape of the acceptance vs. momentum spectrum as a function of the stopping target location.
}
\end{figure}
}

\newcommand{\FigOptimStopTgtPosHeights}{
\begin{figure}[p]
\centering 
\subfloat[][\figlabel{optim:StopTgtPos:Height:42.5}Electrons from 40 to 45 MeV/c]{%
\includegraphics[width=0.95\textwidth,trim=1.15cm 0.05cm 0.3cm 0.180cm,clip]{figs/optimisation/StopTgtPosition/WithBeamBlocker-Height-VaryShifts-Momentum_42-5.pdf}}%
\\\subfloat[][\figlabel{optim:StopTgtPos:Height:62.5}Electrons from 60 to 65 MeV/c]{%
\includegraphics[width=0.95\textwidth,trim=1.15cm 0.05cm 0.3cm 0.180cm,clip]{figs/optimisation/StopTgtPosition/WithBeamBlocker-Height-VaryShifts-Momentum_62-5.pdf}}%
\\\subfloat[][\figlabel{optim:StopTgtPos:Height:82.5}Electrons from 80 to 85 MeV/c]{%
\includegraphics[width=0.95\textwidth,trim=1.15cm 0.05cm 0.3cm 0.180cm,clip]{figs/optimisation/StopTgtPosition/WithBeamBlocker-Height-VaryShifts-Momentum_82-5.pdf}}%
\\\subfloat[][\figlabel{optim:StopTgtPos:Height:102.5}Electrons from 100 to 105 MeV/c]{%
\includegraphics[width=0.95\textwidth,trim=1.15cm 0.05cm 0.3cm 0.180cm,clip]{figs/optimisation/StopTgtPosition/WithBeamBlocker-Height-VaryShifts-Momentum_102-5.pdf}}%
\\\subfloat[][\figlabel{optim:StopTgtPos:Height:102.5-zoom}Electrons from 100 to 105 MeV/c (Zoomed)]{%
\includegraphics[width=0.95\textwidth,trim=1.15cm 0.05cm 0.3cm 0.180cm,clip]{figs/optimisation/StopTgtPosition/WithBeamBlocker-Height-VaryShifts-Momentum_102-5-Zoom.pdf}}%
\caption{\figlabel{optim:StopTgtPos:Height}
The effect of stopping target position on the height of electrons with a fixed momentum as they pass through the Electron Spectrometer.
The size of the variation indicates the stability of the dipole tune; the two parameters are clearly correlated.
Also striking -- particularly in \protect\subref{fig:optim:StopTgtPos:Height:102.5-zoom} -- is the way the dependence on the helical pitch angles is affected by the stopping target position.
}
\end{figure}
}

%\newcommand{\FigOptimStopTgtPosSensitivityIntegral}{
%\begin{figure}[tb]
%\centering 
%%	\fbox{
%\includegraphics[width=0.9\textwidth,trim=1.2cm 0.05cm 0.9cm 0.6cm,clip]{figs/optimisation/StopTgtPosition/Tidied_-AcceptedMomentum-Integrated_adjacent.pdf}
%%}
%\caption{\figlabel{optim:StopTgtPos:AcceptedIntegrated}
%The variation in sensitivity to different momentum electrons as a function of the target position, with respect to the nominal location.
%The darkest red line towards the top of the plot represents the signal sensitivity, and it is that line that should therefore be maximised.
%}
%\end{figure}
%}
%
%\newcommand{\FigOptimStopTgtPosSensitivityIntegralRatio}{
%\begin{figure}[tb]
%\centering 
%%	\fbox{
%\includegraphics[width=0.9\textwidth,trim=1.2cm 0.05cm 0.9cm 0.6cm,clip]{figs/optimisation/StopTgtPosition/Tidied_-AcceptedMomentum-Integrated_ratios.pdf}
%%}
%\caption{\figlabel{optim:StopTgtPos:AcceptedIntegratedRatio}
%The relative avvep
%}
%\end{figure}
%}
